% TOP_PROBLEM: {"id": 1268, "title": "Skolem Problem", "category_id": 1, "category": {"id": 1, "name": "number_theory", "display_name": "Number Theory", "description": "Properties of integers, prime numbers, Diophantine equations.", "slug": "number-theory", "order_index": 1, "created_at": "2026-07-31T15:26:25.670Z"}, "status": "open"}
% FIRST_POSTED: 2026-09-25
% ENTRY_KIND: statement_only
% !TeX program = lualatex
% Definition and sources only; this file is not an LLM solution attempt.
\documentclass[11pt]{article}
\usepackage[margin=25mm]{geometry}
\usepackage{fontspec}
\setmainfont{Latin Modern Roman}
\usepackage{amsmath,amssymb,mathtools}
\usepackage{unicode-math}
\setmathfont{Latin Modern Math}
\usepackage{xurl,hyperref}
\hypersetup{colorlinks=true,urlcolor=blue}
\setcounter{secnumdepth}{0}
\setlength{\parindent}{0pt}
\setlength{\parskip}{5pt}
\setlength{\emergencystretch}{3em}
\begin{document}
\section*{\texorpdfstring{Skolem Problem}{Skolem Problem}}\label{1268}
\subsection{Definitions and mathematical statement}
A linear recurrence sequence of order $k$ is determined by coefficients $c_0,\ldots,c_{k-1}$ and initial terms $u_0,\ldots,u_{k-1}$ through
\[
 u_{n+k}=c_{k-1}u_{n+k-1}+\cdots+c_0u_n\quad(n\ge0).
\]
The data are integers or algebraic numbers given by effective exact encodings. The Skolem problem asks for an algorithm that always halts and decides whether $\exists n\ge0$ with $u_n=0$. The order is part of the input and is unbounded. An algorithm for a fixed low order, or a description of the zero set without effective bounds on the exceptional indices, does not settle the general decision problem.
\par\smallskip\textit{Formulation and concept references:} \href{https://theoretics.episciences.org/17020}{[S1]}.

\subsection{Short English statement}
Can one always decide from a linear recurrence and its initial values whether any term of the resulting sequence is exactly zero?
\subsection{Sources}
\begin{itemize}
\item[S1]Piotr Bacik, "Completing the picture for the Skolem Problem on order-4 linear recurrence sequences," TheoretiCS, Volume 4 (published 2 December 2025), DOI 10.46298/theoretics.25.28; arXiv:2409.01221.
\url{https://theoretics.episciences.org/17020}.
\textit{Formulation source.}
\end{itemize}
\end{document}
